\chapter{Literature Review and Study of Existing Systems}\label{ch:lit}

\section{Study of Existing Systems}
\subsection{The national systems}
India's transport administration runs on a small family of national systems that sit on
top of state-level transport departments. According to the Parivahan portal, MoRTH
entrusted the National Informatics Centre with standardising and deploying
\textbf{Vahan}, used for vehicle registration, and \textbf{Sarathi}, used for driving
licences, and with compiling their data into state and national registers
\cite{parivahan-about}. The citizen-facing front is the \textbf{Parivahan Sewa} portal.
Enforcement uses \textbf{e-challan}, in which traffic penalties are issued and paid
electronically, and the Motor Vehicles (Amendment) Act, 2019 \cite{mva2019} widened the
use of electronic monitoring and revised penalties. States run their own portals and
fee rules on top of these systems.

We studied these systems from their public descriptions only. Their internal designs
are not public, so we make no claim about how they are built. What the public material
does show, and what we used as a requirements source, is the \emph{range of services}:
learner and permanent licences, renewals and duplicates, vehicle registration, transfer
of ownership, fitness, pollution and insurance records, permits, challans and fee
payment.

\subsection{Typical office practice}
Below the national systems, the work done at a counter is still organised around
paper files and shared logins in many offices. This is a general observation, not a
measured one, and we do not attribute it to any particular office. We use it only to
motivate the design questions: who changed a record, what was the previous state, and
can two staff act on the same file at once?

\subsection{Typical student projects}
A common way to build such a system in a course is a small CRUD application with five to
eight tables and rules written only in application code. That approach is simple, but a
second program, a manual SQL fix or a bug can break the rules without anyone noticing.
Our project is deliberately at the other end: the rules are in the schema as well as in
the services.

\begin{table}[H]
\centering\small
\caption{Approaches to RTO record-keeping and what this project takes from each}\label{tab:existing}
\begin{tabularx}{\linewidth}{@{}p{3.2cm}XX@{}}
\toprule
\textbf{Approach} & \textbf{Strength} & \textbf{Limitation for our purpose}\\ \midrule
Paper files and spreadsheets & Flexible, needs no software & No enforced rules, no audit trail, no concurrency control\\
National portals (Vahan, Sarathi, Parivahan Sewa, e-challan) \cite{parivahan-about} & Nationwide scale, uniform services & Internal design not public; cannot be studied or extended by a student team\\
Small CRUD course project & Easy to build and explain & Few tables; rules only in application code; usually untested for concurrency\\
\textbf{This project} & 63 tables, rules in the database and the services, tested concurrency and measured performance & Single-machine prototype; no integration with national registers\\
\bottomrule
\end{tabularx}
\end{table}

\section{Review of Technical Literature and Resources}
\subsection{The relational model and normalisation}
Codd's 1970 paper \cite{codd1970} introduced the relational model: data held as
relations, queried by their content, with integrity kept independent of programs. The
standard textbooks \cite{silberschatz2019,elmasri2016} develop the working tools that
we used: keys, referential integrity, the normal forms and the argument for splitting a
table to remove update anomalies.
\emph{Relevance:} the 63-table design is a direct application; reference data such as
vehicle types and licence classes sit in their own tables, and a person is stored once
while citizens and employees extend that row.
\emph{Limit:} normalisation reduces redundancy but does not decide which rules belong in
the database; that is a design judgement we discuss in Chapter~\ref{ch:design}.

\subsection{Transactions and concurrency}
The idea that a group of operations should succeed or fail as a unit is usually
summarised by the ACID properties, a term from H\"arder and Reuter's survey of
recovery \cite{harder1983}, and treated at length by Gray and Reuter \cite{gray1993}.
Locking reads such as \cmd{SELECT ... FOR UPDATE} are the practical tool: in InnoDB the
statement takes exclusive locks on the rows it reads, other writers wait, and the locks
are released only when the transaction commits or rolls back \cite{mysql-locking}.
\emph{Relevance:} slot booking, ownership transfer, payments and refunds all follow this
pattern, with the added rule that locks are always taken in the same order so that two
transactions cannot wait on each other forever (Chapter~\ref{ch:algo}).
\emph{Limit:} pessimistic locking serialises competing writers; it is the right choice
for hot rows such as the last seat of a slot, but it is not free.

\subsection{Constraints and what a database can enforce}
SQL lets a schema state rules as primary, foreign and unique keys, NOT NULL, CHECK
constraints and generated columns. MySQL 8 enforces CHECK constraints, but its manual
limits what a CHECK expression may contain: ``literals, deterministic built-in
functions, and operators are permitted'', and functions such as \cmd{NOW()} are named as
non-deterministic and therefore not allowed \cite{mysql-check}.
\emph{Relevance:} this rule matters in practice. Two constraints in the supplied script
used the current date and are rejected by the server (Section~\ref{sec:defects}).

\subsection{Role-based access control}
Sandhu et al.\ \cite{sandhu1996} describe role-based access control, in which users are
assigned roles and roles are assigned permissions, so that administrators manage a small
number of roles instead of many user-by-user rules.
\emph{Relevance:} users, roles, permissions and their two junction tables are ordinary
tables in the schema. A permission check is a join, not a hard-coded rule, and the
administrator edits roles from the portal.
\emph{Limit:} RBAC answers ``may this role do this kind of thing'' but not ``may this
user see this particular record''. We added an ownership rule on top: permissions such
as \cmd{application.view_own} restrict a citizen to their own records.

\subsection{Web APIs, tokens and passwords}
Fielding's thesis \cite{fielding2000} defines the REST style: stateless requests, uniform
resource addresses and standard verbs. JSON Web Tokens are specified in RFC~7519
\cite{rfc7519}. Passwords should never be stored in recoverable form; the adaptive
BCrypt scheme of Provos and Mazi\`eres \cite{provos1999} is deliberately slow, and OWASP's
guidance \cite{owasp-pw} recommends such a scheme with a per-password salt.
\emph{Relevance:} the backend is a stateless REST service with short-lived access
tokens and rotating refresh tokens, and passwords are stored as BCrypt hashes
(Section~\ref{sec:security}).

\subsection{Software engineering and algorithm analysis}
The process and documentation practices follow standard texts
\cite{sommerville2015,pressman2019}: a requirements specification with numbered
requirements, UML models, layered design, incremental implementation and testing at
several levels. The analysis of the critical operations in Chapter~\ref{ch:algo}
uses the standard tools of algorithm analysis \cite{cormen2022}, namely asymptotic
cost, ordering arguments, and graph methods such as topological ordering.

\section{Comparative Analysis}
Table~\ref{tab:compare} compares the approaches on the qualities that this project
cares about. The entries for the national systems say ``not publicly documented''
where we could not verify a claim; we do not guess.

\begin{table}[H]
\centering\small
\caption{Comparison on design qualities (``n/p'' means not publicly documented)}\label{tab:compare}
\begin{tabularx}{\linewidth}{@{}p{3.4cm}CCCC@{}}
\toprule
\textbf{Quality} & \textbf{Paper / spreadsheet} & \textbf{National portals} & \textbf{Small CRUD project} & \textbf{This project}\\ \midrule
Rules enforced in the database & No & n/p & Rarely & Yes (CHECK, UNIQUE, generated columns, triggers)\\
Full status history & Manual & n/p & Rarely & Yes (history tables, audit log)\\
Row-level concurrency control & No & n/p & Rarely & Yes (ordered row locks, guarded updates)\\
Permissions stored as data & No & n/p & Rarely & Yes (7 roles, 62 permissions)\\
Idempotent payments & No & n/p & Rarely & Yes (gateway reference)\\
Automated tests published & No & n/p & Rarely & 185 integration tests plus SQL tests\\
Integration with national registers & No & Yes & No & No\\
\bottomrule
\end{tabularx}
\end{table}

\section{Limitations of Existing Approaches}
\begin{enumerate}
\item \textbf{Rules live outside the data.} When rules exist only in forms or application code, any other route to the data can break them.
\item \textbf{History is overwritten.} A single status column shows the present but not the path, which makes audits and appeals harder.
\item \textbf{No concurrency discipline.} Two people acting on the same file at the same moment is the normal case in an office, yet simple systems ignore it.
\item \textbf{Access is coarse.} A few fixed access levels either give staff too much or block them from doing their work, and citizens often cannot be limited to their own records.
\item \textbf{Retries are dangerous.} A repeated payment request can charge twice unless the system recognises it.
\item \textbf{Claims are untested.} Reports often state that a feature works without saying how it was checked.
\end{enumerate}

\section{Identified Engineering Gap}
Existing material covers the individual ideas well: the relational model, transactions,
RBAC, REST and JWT. What is hard to find, at a size that a student team can inspect, is
a worked example that \emph{puts them together on one realistic domain and then proves
the combination holds}, including under concurrent load and with an honest list of the
defects found. That is the gap this project tries to fill.

\section{Proposed Improvements}
We respond to the limitations above as follows:
\begin{itemize}
\item rules are placed in the schema and repeated in the services;
\item every status that matters has a history table, and an append-only audit log records changes with the user and the time;
\item row locks are taken in a fixed order inside transactions, with guarded updates and a database CHECK as the last line of defence;
\item roles and permissions are tables, with an own-records rule for citizens;
\item payments are idempotent by gateway reference, and failed or timed-out attempts are kept as attempts;
\item all of this is covered by automated tests whose results are reported, together with the defects they exposed.
\end{itemize}
